\section{Вывод}

В ходе выполнения лабораторной работы был реализован расширенный калькулятор на платформе Android с использованием языка программирования Kotlin и библиотеки Expression Builder (exp4j). 

Была достигнута поставленная цель --- расширение базового функционала калькулятора дополнительными операторами и функциями согласно варианту 5. Реализованы следующие возможности:

\begin{itemize}
    \item \textbf{Унарные операторы + и - (Sign Operators)}: реализована функция \texttt{normalizeUnaryOperators()}, которая корректно обрабатывает унарные знаки в различных позициях выражения. Калькулятор поддерживает выражения вида \texttt{+2 - (-2)}, где унарный плюс и минус применяются к числам и подвыражениям.
    
    \item \textbf{Функция sinh (гиперболический синус)}: добавлена пользовательская функция через интерфейс \texttt{Function} библиотеки exp4j, использующая стандартную функцию \texttt{kotlin.math.sinh()} для вычисления гиперболического синуса.
    
    \item \textbf{Функция signum (функция знака)}: реализована пользовательская функция, возвращающая знак числа: 1 для положительных чисел, -1 для отрицательных и 0 для нуля.
\end{itemize}

Все функции интегрированы в пользовательский интерфейс через дополнительные кнопки на клавиатуре калькулятора. Реализована обработка ошибок при вычислении выражений с выводом сообщения "Error" в случае некорректного ввода.

Проведено тестирование функционала на следующих примерах:
\begin{itemize}
    \item \texttt{+2 - (-2)} $\rightarrow$ результат: 4
    \item \texttt{sinh(1)} $\rightarrow$ результат: 1.1752011936438014
    \item \texttt{signum(-5)} $\rightarrow$ результат: -1
    \item \texttt{signum(0)} $\rightarrow$ результат: 0
\end{itemize}

Все тесты прошли успешно, функционал работает корректно. Калькулятор готов к использованию и может быть расширен дополнительными функциями в будущем.


